\documentclass[11pt]{article}
\usepackage{fontspec}
\setmainfont{Latin Modern Roman}
\usepackage{amsmath,amssymb,mathtools,unicode-math}
\setmathfont{Latin Modern Math}
\usepackage{xurl,hyperref}
\hypersetup{colorlinks=true,urlcolor=blue,pdftitle={Problem statement}}
\newcommand{\sref}[2]{\hyperlink{source:#1:#2}{[S#2]}}
\newcommand{\cataloguescope}[1]{\paragraph{Mathematical question.}#1}
\begin{document}
\section{Some Questions — Question 47}\label{statement}

\subsection{Definitions and mathematical statement}

\paragraph{Shared notation.}$\mathbb N=\{1,2,\ldots\}$, and $\mathbb Z,\mathbb Q,\mathbb R,\mathbb C$ denote the integers, rationals, reals and complex numbers. Any other conventions are specified locally.


Let $F_n$ be a free group of finite rank $n\geq2$. A pseudocharacter is a homogeneous quasimorphism $q:F_n\to\mathbb R$: its defect
\[D(q)=\sup_{g,h\in F_n}|q(gh)-q(g)-q(h)|
\]
is finite and $q(g^k)=kq(g)$ for all $g\in F_n$ and integers $k$. It is automorphism invariant when $q(\alpha(g))=q(g)$ for every $\alpha\in\operatorname{Aut}(F_n)$ and $g\in F_n$.
\paragraph{Statement recorded in the supplied source.}
Does $F_n$ admit a nonzero automorphism-invariant pseudocharacter for each $n\geq2$? The cited 2023 theorem answers this historical question affirmatively and in fact gives an infinite-dimensional space of such functions. \sref{statement}{2}
\subsection{Short English statement}
Can a nonzero bounded-defect additive function on a finite-rank free group be invariant under every automorphism? The cited theorem gives an affirmative answer.
\subsection{Sources}
\begin{description}
\item[\hypertarget{source:statement:1}{S1}] ULAM.AI and the UnsolvedMath Contributors, UnsolvedMath: A Curated Collection of Open Mathematics Problems, catalogue entry. CC BY 4.0. \url{https://huggingface.co/datasets/ulamai/UnsolvedMath}.
\item[\hypertarget{source:statement:2}{S2}] F. Fournier-Facio and R. D. Wade, \emph{Aut-invariant quasimorphisms on groups}, arXiv:2211.00800v2 (2023), §1, Question A (Abért, 2010, Question 47) and Theorem B. \url{https://arxiv.org/abs/2211.00800}.
\end{description}
\label{end:statement}
\end{document}
